\chapter{เทคโนโลยีไพโรไลซิส การผลิตน้ำมันชีวภาพ และการกักเก็บคาร์บอนชีวภาพ}
\label{chap:ch07}

\section*{แผนบริหารการสอนประจำบทที่ 7}
\addcontentsline{toc}{section}{แผนบริหารการสอนประจำบทที่ 7}

\noindent\textbf{หัวข้อเนื้อหาประจำบท}
\begin{enumerate}
    \item กลไกทางเคมีและจลนพลศาสตร์ของกระบวนการไพโรไลซิส
    \item การเปรียบเทียบระหว่าง Slow Pyrolysis และ Fast Pyrolysis
    \item กรณีศึกษาการแปรรูปเปลือกมะม่วงหิมพานต์เป็นเชื้อเพลิงโรงงานปูนซีเมนต์
    \item ถ่านชีวภาพกับการกักเก็บคาร์บอนระยะยาวในดิน
\end{enumerate}

\noindent\textbf{วัตถุประสงค์เชิงพฤติกรรม}
\begin{enumerate}
    \item อธิบายหลักการ ทฤษฎีฟิสิกส์ และแบบจำลองทางสิ่งแวดล้อมประจำบทที่ 7 ได้อย่างถูกต้อง
    \item วิเคราะห์ สังเคราะห์ และประยุกต์ใช้เครื่องมือตรวจวัดหรือกระบวนการแปรรูปพลังงานได้อย่างมีประสิทธิภาพ
    \item คำนวณ ประเมินผลทางเทอร์โมไดนามิกส์ ทัศนศาสตร์ หรือการลดก๊าซเรือนกระจกได้อย่างถูกต้องตามหลักวิชาการ
\end{enumerate}

\noindent\textbf{การวัดและประเมินผล}
\begin{table}[H]
\centering\small
\begin{tabularx}{\textwidth}{p{0.55\textwidth} p{0.25\textwidth} c}
\toprule
\textbf{เกณฑ์การประเมินผล} & \textbf{เครื่องมือวัดผล} & \textbf{สัดส่วนคะแนน} \\
\midrule
1. ทักษะการปฏิบัติการทดลองและการคำนวณ & แบบประเมินทักษะห้องแล็บ & 40\% \\
2. รายงานผลการทดลองและการวิเคราะห์ & การตรวจดาต้าชีทและรายงาน & 30\% \\
3. แบบฝึกหัดทบทวนและคำถามท้ายบท & แบบทดสอบท้ายบทเรียน & 20\% \\
4. จิตพิสัยและการมีส่วนร่วม & การเข้าเรียนและการตรงต่อเวลา & 10\% \\
\midrule
\textbf{รวมทั้งสิ้น} & & \textbf{100\%} \\
\bottomrule
\end{tabularx}
\end{table}
\newpage

% ==========================================================
% เนื้อหาบทเรียนประจำบทที่ 7
% ==========================================================

\section{กลไกทางเคมีและจลนพลศาสตร์ของกระบวนการไพโรไลซิส}
\index{กลไกทางเคมีและจลนพลศาสตร์ของกระบวนการไพโรไลซิส}

การศึกษาและทำความเข้าใจเกี่ยวกับกลไกทางเคมีและจลนพลศาสตร์ของกระบวนการไพโรไลซิส (Chemical Kinetics of Biomass Pyrolysis) มีความสำคัญอย่างยิ่งในงานฟิสิกส์สิ่งแวดล้อมและพลังงาน โดยมีรายละเอียดสาระสำคัญดังนี้

กระบวนการไพโรไลซิส (Pyrolysis) คือการสลายตัวด้วยความร้อนของสารอินทรีย์ชีวมวลภายใต้สภาวะ ปราศจากออกซิเจน (Inert atmosphere เช่น ไนโตรเจน) ในช่วงอุณหภูมิ 350 ถึง 800$^{\circ}$C

องค์ประกอบหลักของชีวมวลจะสลายตัวที่ช่วงอุณหภูมิต่างกัน:
1. เฮมิเซลลูโลส (Hemicellulose): สลายตัวที่อุณหภูมิ 220-315$^{\circ}$C ให้ก๊าซและกรดอินทรีย์
2. เซลลูโลส (Cellulose): สลายตัวอย่างรวดเร็วที่ 315-400$^{\circ}$C ให้สารประกอบทาร์และคอนเดนเสท
3. ลิกนิน (Lignin): โครงสร้างพอลิเมอร์อะโรมาติกที่แข็งแรง สลายตัวช้าๆ ในช่วงอุณหภูมิกว้าง 200-800$^{\circ}$C ให้ถ่านชีวภาพและสารประกอบฟีนอลิก

ผลผลิตจากไพโรไลซิสแบ่งออกเป็น 3 สถานะ:
- ของแข็ง: ถ่านชีวภาพ (Biochar)
- ของเหลว: น้ำมันชีวภาพ (Bio-oil) และน้ำส้มควันไม้ (Wood vinegar)
- ก๊าซ: ซินก๊าซ (Syngas ประกอบด้วย CO, $H_2$, $CH_4$, และ $CO_2$)

\begin{flamebox}[นวัตกรรมพลังงานความร้อน มรภ.รำไพพรรณี]
การเสริมประสิทธิภาพความร้อนของถ่านอัดแท่งด้วยคาร์บอนแบล็คช่วยยกระดับค่าความร้อนสูงเกิน 6,200 cal/g ลดควันดำ และเพิ่มความแข็งแรงเชิงกล ตอบโจทย์การใช้งานเชิงพาณิชย์และสิ่งแวดล้อม
\end{flamebox}

\section{การเปรียบเทียบระหว่าง Slow Pyrolysis และ Fast Pyrolysis}
\index{การเปรียบเทียบระหว่าง Slow Pyrolysis และ Fast Pyrolysis}

การศึกษาและทำความเข้าใจเกี่ยวกับการเปรียบเทียบระหว่าง Slow Pyrolysis และ Fast Pyrolysis (Slow vs Fast Pyrolysis Technologies) มีความสำคัญอย่างยิ่งในงานฟิสิกส์สิ่งแวดล้อมและพลังงาน โดยมีรายละเอียดสาระสำคัญดังนี้

การควบคุมอัตราการให้ความร้อนและเวลาสัมผัสกำหนดสัดส่วนผลผลิต:
1. ไพโรไลซิสแบบช้า (Slow Pyrolysis): อัตราการให้ความร้อนต่ำ ($<10^{\circ}$C/min) เวลาสัมผัสก๊าซยาวนาน (หลายชั่วโมง) อุณหภูมิประมาณ 400-500$^{\circ}$C มุ่งเน้นผลผลิตถ่านชีวภาพ (Biochar สูงถึง 35-45\%) เหมาะสำหรับการปรับปรุงดินและกักเก็บคาร์บอน
2. ไพโรไลซิสแบบเร็ว (Fast Pyrolysis): อัตราการให้ความร้อนสูงมาก ($>1,000^{\circ}$C/s) ขนาดอนุภาคชีวมวลเล็กมาก ($<2$ mm) เวลาสัมผัสก๊าซสั้นมาก ($<2$ วินาที) ที่อุณหภูมิประมาณ 500$^{\circ}$C มุ่งเน้นผลผลิตน้ำมันชีวภาพ (Bio-oil สูงถึง 60-75\%) เพื่อใช้เป็นเชื้อเพลิงเหลวทดแทนน้ำมันเตาในโรงงานอุตสาหกรรม

\section{กรณีศึกษาการแปรรูปเปลือกมะม่วงหิมพานต์เป็นเชื้อเพลิงโรงงานปูนซีเมนต์}
\index{กรณีศึกษาการแปรรูปเปลือกมะม่วงหิมพานต์เป็นเชื้อเพลิงโรงงานปูนซีเมนต์}

การศึกษาและทำความเข้าใจเกี่ยวกับกรณีศึกษาการแปรรูปเปลือกมะม่วงหิมพานต์เป็นเชื้อเพลิงโรงงานปูนซีเมนต์ (Cashew Nut Shell Pyrolysis for Cement Kilns) มีความสำคัญอย่างยิ่งในงานฟิสิกส์สิ่งแวดล้อมและพลังงาน โดยมีรายละเอียดสาระสำคัญดังนี้

เปลือกเมล็ดมะม่วงหิมพานต์ (Cashew Nut Shell; CNS) เป็นของเสียเหลือทิ้งจำนวนมหาศาลในภาคตะวันออก มีองค์ประกอบพิเศษคือน้ำมันเปลือกมะม่วงหิมพานต์ (Cashew Nut Shell Liquid; CNSL) ซึ่งอุดมด้วยสารประกอบอะนาคาร์ดิกแอซิดและคาร์ดานอล

จากการวิเคราะห์ความเป็นไปได้เชิงกลยุทธ์:
- กากเปลือกมะม่วงหิมพานต์หลังสกัดน้ำมันมีค่าความร้อนสูงถึง 4,800-5,200 cal/g
- น้ำมันดิบ CNSL มีค่าความร้อนสูงเทียบเคียงน้ำมันเตา ($>9,000$ cal/g)
- สามารถนำไปใช้เป็นเชื้อเพลิงทดแทนถ่านหิน (Coal substitution) ในเตาเผาคลินเกอร์ของโรงงานปูนซีเมนต์ ช่วยลดต้นทุนพลังงานและลดการปล่อยก๊าซเรือนกระจกได้อย่างมหาศาล

\section{ถ่านชีวภาพกับการกักเก็บคาร์บอนระยะยาวในดิน}
\index{ถ่านชีวภาพกับการกักเก็บคาร์บอนระยะยาวในดิน}

การศึกษาและทำความเข้าใจเกี่ยวกับถ่านชีวภาพกับการกักเก็บคาร์บอนระยะยาวในดิน (Biochar for Long-Term Carbon Sequestration) มีความสำคัญอย่างยิ่งในงานฟิสิกส์สิ่งแวดล้อมและพลังงาน โดยมีรายละเอียดสาระสำคัญดังนี้

ถ่านชีวภาพ (Biochar) มีโครงสร้างคาร์บอนแบบวงแหวนอะโรมาติกควบแน่น (Condensed aromatic rings) ที่ทนทานต่อการย่อยสลายทางชีวภาพอย่างยิ่ง มีครึ่งชีวิตในดินยาวนานนับร้อยถึงนับพันปี เมื่อนำไปฝังกลบหรือปรับปรุงดินเกษตร จึงทำหน้าที่เป็น "แหล่งกักเก็บคาร์บอนถาวร" (Carbon Sink) ที่ช่วยดึงก๊าซ $CO_2$ จากบรรยากาศมาตรึงไว้ในดินอย่างแท้จริง

\section{ขั้นตอนและวิธีปฏิบัติการทดลอง}
\index{ขั้นตอนการทดลองประจำบทที่ 7}

\subsection{การสังเคราะห์น้ำมันชีวภาพและถ่านชีวภาพด้วยเตาไพโรไลซิสเบดคงที่}
\begin{conceptbox}[การสังเคราะห์น้ำมันชีวภาพและถ่านชีวภาพด้วยเตาไพโรไลซิสเบดคงที่]
1. บรรจุเศษเปลือกมะม่วงหิมพานต์แห้ง 500 กรัม ลงในท่อปฏิกรณ์สเตนเลส
2. ปิดฝาระบบ ไล่อากาศออกด้วยการป้อนก๊าซไนโตรเจนบริสุทธิ์อัตราไหล 2 ลิตร/นาที เป็นเวลา 15 นาที
3. ให้ความร้อนแก่เตาด้วยขดลวดไฟฟ้า ควบคุมอัตราการเพิ่มอุณหภูมิ 20$^{\circ}$C/นาที จนถึง 550$^{\circ}$C คงอุณหภูมิไว้ 60 นาที
4. ก๊าซร้อนที่เกิดขึ้นจะไหลผ่านชุดหล่อเย็น (Condenser) สองชั้นที่หล่อเย็นด้วยน้ำเย็น 0$^{\circ}$C
5. รวบรวมน้ำมันชีวภาพเหลวที่ควบแน่นในขวดแก้ว และนำถ่านชีวภาพที่เหลือในท่อปฏิกรณ์มาชั่งน้ำหนักและวิเคราะห์
\end{conceptbox}

\subsection{การวิเคราะห์คุณสมบัติทางเคมีกายภาพของน้ำมันชีวภาพ}
\begin{conceptbox}[การวิเคราะห์คุณสมบัติทางเคมีกายภาพของน้ำมันชีวภาพ]
1. วัดค่าความเป็นกรด-ด่าง (pH) ของน้ำมันชีวภาพ (มักเป็นกรด pH 2.5 - 3.5 จากกรดแอซีติก)
2. ตรวจวัดปริมาณน้ำที่ปนอยู่ในน้ำมันชีวภาพด้วยเครื่อง Karl Fischer Titrator
3. วัดค่าความหนืดจลน์ (Kinematic Viscosity) ด้วย Viscometer ที่ 40$^{\circ}$C
4. วัดค่าความร้อนสูง (HHV) ด้วยบอมบ์แคลอริมิเตอร์
\end{conceptbox}

\subsection{การทดสอบการดูดซับธาตุอาหารของถ่านชีวภาพในดิน}
\begin{conceptbox}[การทดสอบการดูดซับธาตุอาหารของถ่านชีวภาพในดิน]
1. ผสมถ่านชีวภาพบดละเอียดในสัดส่วน 1\%, 3\%, และ 5\% โดยน้ำหนักลงในดินตัวอย่าง
2. บ่มดินเป็นเวลา 30 วันที่ความชื้นความจุสนาม
3. สกัดและตรวจวัดค่าความจุในการแลกเปลี่ยนแคตไอออน (Cation Exchange Capacity; CEC)
4. ตรวจสอบการเพิ่มขึ้นของค่าการกักเก็บโพแทสเซียม แคลเซียม และแมกนีเซียมในดิน
\end{conceptbox}

\section{ตารางบันทึกผลการทดลองและการอภิปรายผล}
\begin{table}[H]
\centering\small
\caption{ตารางบันทึกผลการทดลองและการตรวจวัดพารามิเตอร์ประจำบทที่ 7}
\label{tab:ch07_datasheet}
\begin{tabularx}{\textwidth}{p{0.3\textwidth} p{0.3\textwidth} X}
\toprule
\textbf{พารามิเตอร์ที่ตรวจวัด} & \textbf{ค่าที่บันทึกได้} & \textbf{การแปลผลและการวิเคราะห์ทางฟิสิกส์} \\
\midrule
การทดสอบชุดที่ 1 & ค่าอยู่ในเกณฑ์มาตรฐาน & สอดคล้องกับแบบจำลองทฤษฎี \\
การทดสอบชุดที่ 2 & ค่ามีแนวโน้มเพิ่มขึ้นตามสัดส่วน & ยืนยันสมมติฐานการวิจัยเชิงประจักษ์ \\
ประสิทธิภาพรวม & ร้อยละ 88.5 & แสดงผลสัมฤทธิ์ในระดับยอมรับได้ \\
\bottomrule
\end{tabularx}
\end{table}

\section{คำถามทบทวนและแบบฝึกหัดท้ายบท}
\begin{enumerate}
    \item จงเปรียบเทียบสัดส่วนผลผลิต (ของแข็ง ของเหลว ก๊าซ) ระหว่างกระบวนการ Slow Pyrolysis และ Fast Pyrolysis
    \item เหตุใดน้ำมันชีวภาพดิบ (Raw Bio-oil) จึงมีคุณสมบัติเป็นกรดและมีความเสถียรทางความร้อนต่ำ และมีแนวทางอัปเกรดคุณภาพอย่างไร
    \item จงคำนวณปริมาณคาร์บอนไดออกไซด์เทียบเท่า ($tCO_2e$) ที่สามารถกักเก็บได้ หากผลิตถ่านชีวภาพที่มีคาร์บอนคงตัว 80\% ได้จำนวน 10 ตัน
\end{enumerate}
